\RequirePackage[symbols,acronym,xindy={language=german-din},toc,nostyles]{glossaries}
\RequirePackage{tabularx}
\RequirePackage{ltablex}  % ltablex kombiniert longtable und tabularx, optional
\RequirePackage{booktabs}
\RequirePackage{environ}

\glsdisablehyper

\keepXColumns  % für tabularx in mehrspaltigen Umgebungen, gut so

\makeglossaries

% =============================
% Definition der Symbolgruppen
% =============================

% Formelzeichen
\newglossaryentry{Formelzeichen}{%
  type=symbols,%
  name={\bfseries Formelzeichen},%
  description={\nopostdesc},%
  sort={a}%
}

% Griechische Buchstaben
\newglossaryentry{greekletter}{%
  type=symbols,%
  name={\bfseries Griechische Buchstaben},%
  description={\nopostdesc},%
  sort={b}%
}

% Indices
\newglossaryentry{Indices}{%
  type=symbols,%
  name={\bfseries Indizierungen},%
  description={\nopostdesc},%
  sort={c}%
}

% ===================================
% Neue Symbole per Makros definieren
% ===================================

\newcommand*{\newformulasymbol}[4][]{%
  \newglossaryentry{#2}{%
    type=symbols,
    name={#2},
    description={\nopostdesc},
    symbol={\ensuremath{#3}},
    user1={\ensuremath{\left[\mathrm{#4}\right]}},
    sort={#3},
    parent={Formelzeichen}
  }%
}

\newcommand*{\newgreek}[4][]{%
  \newglossaryentry{#2}{%
    type=symbols,
    name={#2},
    description={\nopostdesc},
    symbol={\ensuremath{#3}},
    user1={\ensuremath{\left[\mathrm{#4}\right]}},
    sort={#3},
    parent={greekletter}
  }%
}

\newcommand*{\newindice}[3][]{%
  \newglossaryentry{#2}{%
    type=symbols,
    name={#2},
    description={\nopostdesc},
    symbol={\ensuremath{()_{\text{#3}}}},
    user1={\ensuremath{\left[\mathrm{-}\right]}},
    sort={#3},
    parent={Indices}
  }%
}

% ===== Anzeigeformat für Symbole =====
\defglsentryfmt[symbols]{%
  \ifmmode
    \glssymbol{\glslabel}%
  \else
    \glsgenentryfmt~\glsentrysymbol{\glslabel}%
  \fi
}

% ===== Glossar-Stil für Symbolverzeichnis als Tabelle =====
\NewEnviron{tableenv}{%
  \renewcommand{\arraystretch}{1.2}%
  \begin{tabularx}{\columnwidth}{l@{\hspace{4em}} l@{\hspace{4em}} X}%
    \BODY
  \end{tabularx}%
}

\newglossarystyle{symboltablr}{%
  \newcommand*\symbollevel{-1}%
  \renewenvironment{theglossary}{\tableenv}{\endtableenv}%
  \renewcommand*{\glsgroupheading}[1]{}%
  \renewcommand*{\glsgroupskip}{}%
  \newcommand*\symbolhead{%
    \toprule
    \bfseries Symbol & \bfseries Einheit & \bfseries Bedeutung \\ \midrule
  }%
  \renewcommand*{\glossaryheader}{%
    \endfirsthead
    \symbolhead
    \endhead
    \midrule
    \endfoot
    \bottomrule
    \endlastfoot
    \gdef\symbollevel{-1}%
  }%
  \renewcommand*{\glossentry}[2]{%
    \ifglshaschildren{##1}{%
      \ifnum\symbollevel>0
        \tabularnewline\bottomrule\tabularnewline[\smallskipamount]%
      \fi
      \gdef\symbollevel{0}%
      \tabularnewline[\arraystretch\dimexpr-\ht\strutbox-\dp\strutbox\relax]%
      \multicolumn{3}{l}{\glossentryname{##1}}%
    }%
  }%
  \renewcommand*{\subglossentry}[3]{%
    \ifnum\symbollevel=0
      \\[\medskipamount]\symbolhead
    \fi
    \gdef\symbollevel{##1}%
    \glstarget{##2}{\glossentrysymbol{##2}} & \glsentryuseri{##2} & \glossentryname{##2} \\
  }%
}

% ===== Glossar-Stil für Abkürzungsverzeichnis als Tabelle =====
\NewEnviron{acroenv}{%
  \renewcommand{\arraystretch}{1.2}%
  \begin{tabularx}{\columnwidth}{l@{\hspace{4em}} X}%
    \BODY
  \end{tabularx}%
}

\newglossarystyle{acrotab}{%
  \renewenvironment{theglossary}{\acroenv}{\endacroenv}%
  \renewcommand*{\glsgroupheading}[1]{}%
  \renewcommand*{\glsgroupskip}{}%
  \newcommand*\acrohead{%
    \toprule
    \bfseries Abkürzung & \bfseries Bedeutung \\ \midrule
  }%
  \renewcommand*{\glossaryheader}{%
    \acrohead
    \endfirsthead
    \acrohead
    \endhead
    \midrule
    \endfoot
    \bottomrule
    \endlastfoot
  }%
  \renewcommand*{\glossentry}[2]{%
    \glsentryitem{##1}%
    \glstarget{##1}{\glossentryname{##1}} & \glsentrydesc{##1} \\
  }%
}

% ===== Alternative Darstellung Abkürzungsverzeichnis (Liste) =====
\newglossarystyle{acrolist}{%
  \renewenvironment{theglossary}{\begin{description}}{\end{description}}%
  \renewcommand*{\glsgroupheading}[1]{}%
  \renewcommand*{\glsgroupskip}{}%
  \renewcommand*{\glossentry}[2]{%
    \item[\glsentryitem{##1}\glstarget{##1}{\glossentryname{##1}}] \glsentrydesc{##1}%
  }%
}
